\section{Experimental Design}

\subsection{\routea{}: real mechanism rounds}

The campaign begins from an immutable negative result: existing sparse
equation-discovery methods degraded under observation noise. Sparse
identification of nonlinear dynamics provides the underlying formulation
\cite{brunton2016sindy}, and PySINDy supplies a reference implementation
pattern \cite{desilva2020pysindy}. Operon is included as a strong genetic
programming baseline \cite{burlacu2020operon}. The benchmark source contains
ODE and PDE systems under multiple noise settings
\cite{ziaei2025mdbench}. None of these citations is treated as evidence for
our candidate's success; success is determined only by the frozen local
results.

Before the current campaign, an archive audit excludes every system in two
historical panels. It finds enough unused ODE systems to freeze two mutually
exclusive unseen panels of six systems. Each round uses two development
systems, one six-system unseen panel, clean and SNR20 conditions, and three
fresh seeds. The candidate, Operon baseline, and three mechanism-specific
ablations form 240 evaluated cells per round. The unseen result is revealed
only if the development median improvement is at least 0.15 and all required
candidate and baseline cells are valid.

Round 1, \path{noise_conditioned_ensemble_sindy}, applies
noise-conditioned Savitzky--Golay derivative estimation and coefficient
bootstrap aggregation. Its ablations remove noise conditioning, ensemble
aggregation, or smoothing. Round 2,
\path{spline_group_sparse_sindy}, replaces the derivative path with
analytic smoothing-spline derivatives and projects coefficients through
cross-output group sparsity. Its ablations remove spline derivatives, group
projection, or adaptive thresholding. Round 2 is linked to the Round 1
negative-result hash, so the change is a successor mechanism rather than an
independent retrospective comparison.

The preregistered method gate requires a system-level, failure-aware bootstrap
95\% confidence interval whose lower endpoint exceeds zero, three-seed
reproduction, the strong baseline, all ablations, and an idempotent rerun.
Failure on any item produces a negative report. No threshold is weakened and
no revealed panel is exchanged.

\subsection{\routeb{}: controlled research-workflow matrix}

After both \routea{} rounds fail their unseen gate, the campaign freezes the
systems study. The sources comprise \UciTaskCount{} fresh local executions
over UCI datasets \cite{dua2019uci} and \MDBenchTaskCount{} replayed trace
summaries from the revealed equation-discovery campaign. Each source record
contains file hashes, validated scalar metrics, a positive or negative truth
label, and an explicit use restriction. UCI runs establish that the workflow
ingests newly computed real evidence. MDBench traces establish continuity
with the failed research rounds. Because those traces are already revealed,
they test controller behavior only.

We inject one frozen workflow state into each otherwise valid source task.
Table~\ref{tab:tasks} shows the matrix. Faults cover source integrity,
metric direction, configuration freeze, missing evidence maps, unsupported
claims, missing seed coverage, failed experiment cells, and
claim-evidence mismatch. Three tasks require a prior \vault{} record for
repair. Two faults are visible during planning; the remainder require
execution feedback or the evidence gate. One task in each family has no
fault, preventing a controller from benefiting simply by always changing the
input.

\begin{table*}[t]
\centering
\caption{Frozen task and controlled-fault matrix. Source evidence is real;
the workflow fault is deliberately introduced to isolate controller
behavior. MDBench rows reuse revealed traces and are not method holdouts.}
\label{tab:tasks}
\small
\begin{tabular}{llll}
\toprule
Family & Source task & Initial workflow state & Required repair \\
\midrule
UCI & PenDigits prototypes & none & retain valid evidence \\
UCI & Letter prototypes & stale source hash & rehash source \\
UCI & Spambase prototypes & wrong metric direction & bind direction \\
UCI & Skin prototypes & missing evidence map & restore from \vault{} \\
\midrule
MDBench & binocular rivalry & none & retain valid evidence \\
MDBench & interacting magnets & unfrozen configuration & freeze configuration \\
MDBench & oscillator death & unsupported claim & evidence-bound decision \\
MDBench & population growth & incomplete seed coverage & resume checkpoint \\
MDBench & RC circuit & failed experiment cell & classify and recover \\
MDBench & Van der Pol & claim-evidence mismatch & rewrite claim \\
\bottomrule
\end{tabular}
\end{table*}

\subsection{Controllers and ablations}

The \oneshot{} controller emits one result from the initial state. The
\executeonce{} controller plans once and repairs only a fault that is visible
at planning time; it has no post-execution recovery. The \systemloop{}
preregisters, executes, classifies visible failures, obtains the allowed
\vault{} context, performs the task's frozen repair action, reruns once, and
enforces claim support. These are deterministic transition policies in the
evaluation harness. A local Qwen call provides a structured policy summary
for each main mode, but the summary cannot change transitions or scoring.

Four ablations remove one contract:

\begin{itemize}
  \item \textbf{No \vault{} memory} blocks repairs whose evidence card exists
  only in persistent project memory.
  \item \textbf{No failure feedback} prevents the second attempt after a
  runtime failure.
  \item \textbf{No preregistration} makes the full-loop completion contract
  false for every evaluated cell.
  \item \textbf{No evidence gate} hides claim-support failures from recovery,
  allowing an unsupported claim to survive.
\end{itemize}

Each of seven modes runs on \TaskCount{} tasks and \SeedCount{} seeds,
producing \MainCellCount{} main cells and \AblationCellCount{} ablation cells.
The total is \CellCount{}. An exact reproduction recomputes the scientific
result hash for every cell. Seeds are included in each cell identity even
though the present deterministic transition harness does not sample a
stochastic repair policy. We retain this redundancy to detect accidental
seed-dependent behavior and to keep the protocol compatible with future
stochastic controllers.

\subsection{Metrics and statistical gate}

Task success means that source integrity, experiment completion, evidence-map
completion, reproduction readiness, claim support, and, for full-loop modes,
preregistration are all true after the allowed actions. Negative-result
recovery is counted only when an initially failed cell makes a second attempt
and then succeeds. An unsupported claim is a report claim whose polarity or
scope does not match the frozen source record. Research-decision intervention
counts any human selection after campaign start. External cost records paid
API or cloud-compute usage.

For task \(t\) and seed \(s\), the paired difference is
\[
d_{t,s} =
\mathbb{1}[\text{full loop succeeds}]
-\mathbb{1}[\text{execute-once succeeds}].
\]
We report the mean over the 30 task-seed pairs. A percentile bootstrap samples
the paired differences with replacement for \BootstrapResamples{} resamples
using a frozen random seed, following the nonparametric bootstrap
\cite{efron1979bootstrap}. The internal systems gate passes only if the
confidence interval lower endpoint is above zero, exact reproduction is at
least 0.90, the full loop emits no unsupported claim, intervention is zero,
and every ablation cell completes.
